\section{Un sistema simple y complejo a la vez: código abierto, muro de datos y decisiones de arquitectura}

La sección anterior trató los objetivos de diseño de K2; esta examina las decisiones a nivel de sistema. ¿Por qué pasar de código cerrado a código abierto? ¿Por qué ya es un buen resultado que la multimodalidad no dañe el «cerebro»? ¿Por qué, cuando la Scaling Law se topa con el muro de datos, hay que reequilibrar los FLOPs, el RL, la retroalimentación y la eficiencia de la arquitectura? Estas preguntas muestran que una empresa de modelos no ajusta una sola perilla, sino que toma decisiones continuamente dentro de un sistema complejo.

\begin{figure}[H]
\centering
\includegraphics[width=0.94\textwidth]{open-source-strategy.png}
\caption{De código cerrado a código abierto: el grado de madurez del modelo y del producto, la difusión en el ecosistema y la vía comercial influyen en conjunto en la decisión de abrir el código. Diagrama conceptual propio, elaborado a partir de la conversación entre 00:54:08 y 01:05:00.}
\end{figure}

\begin{knowledgebox}{Lectura de la figura: abrir el código es una decisión técnica y también una decisión de ecosistema}
El código cerrado facilita controlar la experiencia del producto y la comercialización; el código abierto favorece la difusión en el ecosistema, la verificación por parte de los desarrolladores y la construcción de confianza. Que K2 sea de código abierto implica que debe demostrar la capacidad del modelo base y, al mismo tiempo, permitir que el ecosistema Agentic crezca a su alrededor.
\end{knowledgebox}

\subsection{Scaling Law, muro de datos y FLOPs}

La sección anterior trató el código abierto y el ecosistema; esta vuelve a las restricciones duras del entrenamiento de modelos. Yang Zhilin (杨植麟) considera que la Scaling Law se ha topado con el muro de datos, y que es un hecho objetivo. Cuando el crecimiento de los datos de alta calidad se desacelera, seguir haciendo scale de los datos ya no es tan sencillo como antes; por eso aumenta la importancia del FLOPs scaling, del RL, del feedback del entorno, del volante de datos y de la eficiencia de la arquitectura. Por ahora, el scaling basado en FLOPs sigue funcionando, pero el punto de equilibrio podría cambiar en el futuro.

\begin{figure}[H]
\centering
\includegraphics[width=0.96\textwidth]{scaling-data-wall.png}
\caption{Scaling y muro de datos: después del muro de datos, se reequilibran los FLOPs, el RL, la retroalimentación y la eficiencia de la arquitectura. Diagrama conceptual propio, elaborado a partir de la conversación entre 01:03:00 y 01:10:00.}
\end{figure}

\begin{warningbox}{El muro de datos no significa que el entrenamiento se detenga}
El muro de datos no quiere decir que ya no haya datos, sino que escasean los datos de alta calidad, con poco ruido y capaces de aportar una mejora marginal de capacidad. Las empresas de modelos todavía pueden avanzar mediante mejores proporciones de datos, datos sintéticos, RL, entornos de herramientas y test-time scaling.
\end{warningbox}

\begin{knowledgebox}{Cuatro caminos después del muro de datos}
\begin{tabularx}{\textwidth}{p{0.20\textwidth}p{0.34\textwidth}X}
\toprule
Ruta & Qué resuelve & Costo \\
\midrule
FLOPs scaling & Seguir bajando la loss con más cómputo & Costo alto; depende del suministro de hardware. \\
Data efficiency & Obtener más capacidad con los mismos datos & Requiere mejor filtrado, reescritura, proporciones y optimizadores. \\
RL/feedback & Crear nuevas señales de aprendizaje a partir de la retroalimentación del entorno & El ruido del reward y los problemas de seguridad son más difíciles. \\
Architecture & Reducir el costo de contexto largo e inferencia con nuevas estructuras & Puede introducir bias en la capacidad. \\
\bottomrule
\end{tabularx}
\end{knowledgebox}

\subsection{Linear Attention y el riesgo para el «coeficiente intelectual»}

Después del muro de datos, la eficiencia de la arquitectura se vuelve más atractiva; esta sección analiza el riesgo de ese atractivo. En la entrevista se menciona que muchas arquitecturas de long context afectan el «coeficiente intelectual»; una Linear Attention pura podría afectar la capacidad del modelo debido al bias de la arquitectura. Este juicio coincide con la revisión de arquitecturas de Attention del EP119: reducir el costo del contexto largo es importante, pero si el cambio de arquitectura sacrifica capacidad de razonamiento, de memoria o expresiva, puede afectar el techo del modelo.

\begin{figure}[H]
\centering
\includegraphics[width=0.94\textwidth]{linear-attention-iq-risk.png}
\caption{Arquitecturas de Long Context y riesgo para el coeficiente intelectual: la Linear Attention reduce costos, pero el bias puede afectar la capacidad del modelo. Diagrama conceptual propio, elaborado a partir de la conversación entre 01:17:17 y 01:20:00.}
\end{figure}

\begin{knowledgebox}{Términos clave: decisiones de arquitectura}
\begin{tabularx}{\textwidth}{p{0.22\textwidth}p{0.32\textwidth}X}
\toprule
Ruta & Beneficio & Riesgo \\
\midrule
Full Attention & Gran capacidad expresiva; experiencia consolidada & Costo alto en contexto largo; mucha presión sobre la KV cache. \\
Linear Attention & Costo bajo; adecuada para secuencias largas & El bias de la arquitectura puede afectar la capacidad. \\
Sparse Attention & Reduce el cómputo y conserva parte de la visión global & El mecanismo de selección y la eficiencia en hardware son complejos. \\
Hybrid & Punto intermedio entre eficiencia y capacidad & Proporciones, estabilidad del entrenamiento y evaluación más complejas. \\
\bottomrule
\end{tabularx}
\end{knowledgebox}

\subsection{La frontera entre las empresas de modelos base y las empresas de productos Agent}

Lo anterior trató las decisiones de arquitectura internas del modelo; esta sección lleva la pregunta a la frontera de la industria. A largo plazo, la frontera entre las empresas de modelos base y las empresas de aplicaciones Agent no es fija. Las empresas de modelos base controlan la capacidad del modelo, los datos de entrenamiento, la API y el ecosistema de plataforma; las empresas de aplicaciones Agent controlan los escenarios, los flujos de trabajo, el punto de acceso de los usuarios y los datos de retroalimentación. Quien logre convertir la capacidad del modelo en un ciclo cerrado de tareas reales ocupará la posición de mayor valor.

\begin{figure}[H]
\centering
\includegraphics[width=0.94\textwidth]{agent-product-boundary.png}
\caption{Empresas de modelos base vs. empresas de aplicaciones Agent: la frontera a largo plazo depende del modelo, los datos, las herramientas y el punto de acceso de los usuarios. Diagrama conceptual propio, elaborado a partir de la conversación entre 01:09:00 y 01:15:00.}
\end{figure}

\begin{importantbox}{Criterio sobre la frontera}
Si un producto Agent es solo una capa delgada de workflow, la empresa de modelos base absorberá esa frontera; si el producto Agent cuenta con escenarios sólidos, un fuerte retorno de datos y un punto de acceso de usuarios de uso frecuente, puede, por el contrario, definir las necesidades del modelo.
\end{importantbox}

\begin{knowledgebox}{Comparación de modelos de negocio: API, producto y ecosistema}
\begin{tabularx}{\textwidth}{p{0.20\textwidth}p{0.34\textwidth}X}
\toprule
Modelo & Ventaja & Riesgo \\
\midrule
API & Fácil de escalar; los desarrolladores pueden integrarse & Competencia en el precio por token; lejos del valor para el usuario. \\
Producto Agent & Más cerca de las tareas y del punto de acceso de los usuarios & Requiere producto, herramientas, retroalimentación y fiabilidad. \\
Ecosistema de código abierto & Difusión rápida; gran confianza y participación de desarrolladores & Más difícil capturar valor comercial; costo de soporte alto. \\
Producto de código cerrado & Experiencia controlable; ciclo comercial claro & Difusión en el ecosistema y verificación externa limitadas. \\
\bottomrule
\end{tabularx}
\end{knowledgebox}

El modelo de negocio tiene además un problema del lado de los costos: las tareas Agentic suelen consumir más cómputo en inferencia que una conversación común. Una pregunta y respuesta sencilla puede requerir una sola llamada al modelo; una tarea Agent compleja puede requerir planificación, búsqueda, llamadas a herramientas, verificación, reintentos y resumen. Esto amplía la brecha entre el «valor para el usuario» y el «costo de inferencia»: si la tarea es de alto valor, el costo del Agent puede absorberse; si la tarea es solo una conversación de bajo valor, el costo puede aplastar fácilmente el margen bruto.

\begin{knowledgebox}{Estructura de costos de un Agent}
\begin{tabularx}{\textwidth}{p{0.20\textwidth}p{0.34\textwidth}X}
\toprule
Rubro de costo & De dónde proviene & Cómo reducirlo \\
\midrule
Llamadas al modelo & Planificación, ejecución y reflexión en varias rondas & Enrutamiento a modelos pequeños, caché, división de tareas. \\
Llamadas a herramientas & Búsqueda, código, navegador, API & Mejores protocolos de herramientas y recuperación ante fallos. \\
Costo de verificación & Pruebas, revisiones, confirmación humana & Verificadores automáticos y clasificación por nivel de riesgo. \\
Costo de contexto & Historial largo, archivos, bases de código, páginas web & Compresión de memoria y estrategias de recuperación. \\
\bottomrule
\end{tabularx}
\end{knowledgebox}

\subsection{Resumen de la sección}

Las decisiones de sistema detrás de K2 abarcan la estrategia de código abierto, el muro de datos, el FLOPs scaling, la eficiencia de la arquitectura y la frontera del producto. No es una competencia técnica en un solo punto, sino una optimización con múltiples variables que la empresa de modelos realiza dentro de un sistema complejo.
